\section{Introduction}\label{sec:intro}

Haiden, Katzarkov, Kontsevich and Pandit \cite{hkkp26} propose a notion of
K\"ahler geometry on a triangulated category with a stability condition in
the sense of Bridgeland \cite{bri07}, in which an object carries a space of
metrics, an endomorphism $\Omega$ depending on the metric whose trace is the
central charge, a phase $\arg\Omega$ defined by a positivity condition, a
mass $\tr|\Omega|$, and a flow $h^{-1}\dot h=-\arg\Omega$ of the metrics,
called the minimizing flow in \cite[\S4.1]{hkkp26}. A metric is
\emph{harmonic} if $\arg\Omega$ is a scalar, and harmonic metrics are
expected to exist on polystable objects and to be the limits of the flow
\cite[\S\S1,~4.1]{hkkp26}, a categorical form of the theorem of Donaldson,
Uhlenbeck and Yau \cite{don85,uy86}. This paper proves that on a quiver
without oriented cycles the flow converges, up to a scalar function of the
time, on every polystable representation, and that for complexes of
representations with one differential the metrics induced on the cohomology
converge in the same sense when the cohomology is polystable. When it is not
polystable, the long-time behaviour is described through the
Harder--Narasimhan filtration of the cohomology and, on its factors, through
gradient flows of Kempf--Ness functions.

\subsection{Representations}
Let $Q$ be a finite quiver without oriented cycles, with central charges
$z_i$ in the upper half-plane at its vertices, and let $\mathsf R$ be a
representation of $Q$ with Hermitian metrics $h_i$ on its spaces. Write
$\beta^\dagger$ for the adjoint of a linear map for the metrics. The
\emph{moment map} at the vertex $i$ is the Hermitian endomorphism $\mu_i$
equal to the sum of $\beta^\dagger\beta$ over the arrows $\beta$ out of $i$
minus the sum of $\beta\beta^\dagger$ over the arrows into $i$. For a
representation the endomorphism $\Omega$ of \cite[\S5.1]{hkkp26} is
$\bigoplus_i(z_i+\mu_i)$; it is normal, and the flow is
\[
 h_i^{-1}\dot h_i=-\arg(z_i+\mu_i),
\]
with the argument taken by the functional calculus, in $(0,\pi)$. Let
$Z(\mathsf R)=\sum_iz_i\dim\mathsf R_i=|Z|e^{\ii\varphi}$ be the central
charge of $\mathsf R$, and
\[
 \theta_i=\frac{\Ima(e^{-\ii\varphi}z_i)}{\sin\varphi}.
\]
A representation with central charge of argument $\varphi$ is
\emph{semistable} if no subrepresentation has a central charge of argument
larger than $\varphi$; this is King's stability \cite{kin94} for the real
parameter $-\theta$, and a metric is harmonic exactly when
$\mu_i=\theta_i$ for all $i$, that is, when it is a zero of the moment map
shifted by the stability parameter (Section~\ref{sec:convergence}).

\begin{theorem}[Representations]\label{thm:introR}
For all charges in the upper half-plane, every nonzero representation
$\mathsf R$ of $Q$ and every initial metric, along the flow:
\begin{enumerate}
\item[(i)] the Kempf--Ness function
 $\sum_a\norm{\rho_a}^2+\sum_i\theta_i\log\det h_i$ of $\mathsf R$ does not
 increase;
\item[(ii)] if $\mathsf R$ is semistable, then $\mu_i(t)\to\theta_i$ for all
 $i$, and the mass tends to $|Z|$;
\item[(iii)] if $\mathsf R$ is polystable, the metrics $e^{\varphi t}h(t)$
 converge exponentially to a harmonic metric.
\end{enumerate}
\end{theorem}

Here $\rho_a$ is the map of the arrow $a$ and $\norm\cdot$ the
Hilbert--Schmidt norm for the metrics. The factor $e^{\varphi t}$ is needed
because a harmonic metric is not stationary: it moves by the scalar
$e^{-\varphi t}$. After multiplication of the metrics by $e^{\varphi t}$ the
velocity of the flow at the vertex $i$ is $\varphi-\arg(z_i+\mu_i)$, an
increasing function of $\mu_i-\theta_i$ that vanishes at zero. The downward
gradient flow of the Kempf--Ness function, which \cite[(5.5)]{hkkp26}
describes as a simplification of this flow, has the velocity
$\mu_i-\theta_i$ up to a positive factor; its exponential convergence on
polystable representations is classical (see \cite[\S5.1]{hkkp23}, and
\cite{grs21} for the general theory), and the gradient flow of the norm
square of the moment map on the space of representations is studied in
\cite{hw11}. Theorem~\ref{thm:introR} concerns the flow $-\arg\Omega$ itself.
The precise statement, with the rate, is Theorem~\ref{thm:representations}.

\subsection{Complexes with one differential}
For complexes of representations $\Omega$ is defined through a
superconnection \cite[\S5.4]{hkkp26}; it is then not normal in general, and
we take $\arg\Omega$ to be the argument of the unitary factor in the polar
decomposition, which is our reading of the positivity condition of
\cite[\S4.2]{hkkp26}; it agrees with the functional calculus of
\cite[\S5.4]{hkkp26} when $\Omega$ is normal (Section~\ref{sec:setting}). All
statements below refer to this choice, and we call the flow the \emph{phase
flow}.

Fix a vertex $i_0$. A complex of representations $\mathsf E$ of $Q$ has
\emph{a surjective differential at $i_0$} if at each vertex $i\neq i_0$ it is
a vector space in degree zero and at $i_0$ it is $V\xrightarrow{D}U$ in
degrees $0,1$ with $D$ surjective (Definition~\ref{def:type}). Its cohomology
is a representation $\mathcal H$ of $Q$ in degree zero, with $\ker D$ at
$i_0$; we assume $\mathcal H\neq0$, and put
$Z=Z(\mathcal H)=|Z|e^{\ii\varphi}$ and $\theta$ as above. For Hermitian
metrics $h$ on the spaces of $\mathsf E$, the endomorphism $\Omega$ preserves
each of these spaces. We write $\Omega_V$, $\Omega_U$ and $\Omega_i$
($i\neq i_0$) for its blocks,
$\mass(h)=\tr|\Omega_V|+\tr|\Omega_U|+\sum_{i\neq i_0}\tr|\Omega_i|$ for the
mass, where $|T|=(TT^\dagger)^{1/2}$, and $\mu_i$ for the moment map of
$\mathcal H$ with the induced metrics, so that the endomorphism $\Omega$ of
the representation $\mathcal H$ is $\bigoplus_i(z_i+\mu_i)$. Let $\lambda$ be
the least eigenvalue of $DD^\dagger$ and $\norm\cdot_1$ the trace norm. The
quotient complex $\mathsf E/\mathcal H$ is supported at $i_0$ and is exact.
We call its metric realization $(\ker D)^\perp\to U$ the \emph{exact part}.

\begin{theorem}[Convergence on the cohomology]\label{thm:introB}
Let $\mathsf E$ be a complex with a surjective differential at $i_0$. For all
charges in the upper half-plane and every initial metric, along the phase
flow:
\begin{enumerate}
\item[(i)] if $\mathcal H$ is semistable, then $\mu_i(t)\to\theta_i$ for all
 $i$ and $\mass(h(t))\to|Z|$;
\item[(ii)] if $\mathcal H$ is polystable, the metrics induced on the spaces
 of $\mathcal H$, multiplied by $e^{\varphi t}$, converge exponentially to a
 harmonic metric of $\mathcal H$.
\end{enumerate}
\end{theorem}

\begin{theorem}[Decay of the exact part]\label{thm:introA}
Let $\mathsf E$ be a complex with a surjective differential at $i_0$ and
$U\neq0$. For all charges in the upper half-plane and every initial metric,
along the phase flow, $\lambda(t)\to\infty$ and, for $t\geq t_1$,
\[
 \norm{\Omega_V-(z_{i_0}+\mu_{i_0})\oplus0}_1+\norm{\Omega_U}_1
 +\sum_{i\neq i_0}\norm{\Omega_i-(z_i+\mu_i)}_1
 \leq C\,e^{-\varkappa(t-t_1)},
\]
where $\oplus\,0$ refers to $V=\ker D\oplus(\ker D)^\perp$, the constants
$C,\varkappa>0$ depend only on the quiver, the vertex $i_0$, the charges and
the dimensions, and $t_1$ depends in addition on the initial metric.
\end{theorem}

Theorem~\ref{thm:introA} makes no assumption on $\mathcal H$. For $U\neq0$
the complex itself carries no harmonic metric, since the phase has its
spectrum in $(0,\pi)$ on $V$ and in $(-\pi,0)$ on $U$; what is proved is that
the exact part decays and that the metrics induced on the cohomology
converge. By duality \cite[Theorem~5.1]{wue26c} the two theorems hold for
complexes \emph{with an injective differential at $i_0$}, which at $i_0$ are
$U'\to V'$ in degrees $-1,0$ with an injective differential
(Proposition~\ref{prop:injective}). The precise statements, with the
constants, are Theorems~\ref{thm:decay} and~\ref{thm:convergence},
Proposition~\ref{prop:sink} and Remark~\ref{rem:vanishing}.

\subsection{The method}
For a complex with $U\neq0$ the Kempf--Ness function of the cohomology
decreases only up to an error controlled by the exact part
(Theorem~\ref{thm:convergence}(i)), so the proof cannot start from a
Lyapunov function; it rests on an a priori bound for the arrows. Once the arrows are
bounded, the exact part decays, the error becomes exponentially small, and
the convexity of the Kempf--Ness function gives the convergence
(Proposition~\ref{prop:perturbed}). For an arrow $a$ with map $\rho_a$ let
$e_a=\norm{\rho_a}^2$ if $a$ does not start at $i_0$ and
$e_a=\tr(\rho_aR(D^\dagger D)\rho_a^\dagger)$, with $R(b)=(1-e^{-b})/b$, if it
does, and put $\cN=\sum_ae_a$. For a representation $\cN$ is the sum of the
squared norms of the arrows.

\begin{theorem}[A priori bound]\label{thm:introC}
Let $\mathsf E$ be a complex with a surjective differential at $i_0$. There
are constants $\gamma>0$ and $c\geq0$, depending only on the quiver, the
vertex $i_0$, the charges and the dimensions, such that for all metrics
\[
 \frac{d\cN}{dt}\leq-\gamma\,\cN+c
\]
along the phase flow. For a representation one can take $\gamma=\pi/\ell$,
with $\ell$ the number of arrows of a longest path in $Q$, taking $\ell=1$
when $Q$ has no arrows.
\end{theorem}

The derivative $d\cN/dt$ is a sum of three traces, one for $V$, one for $U$
and one for the remaining vertices (Lemma~\ref{lem:energyderivative}). The
trace for $V$ is bounded by a trace inequality for the polar phase, and the
trace for $U$ by a bound for the negative part of a Hermitian matrix that
occurs in $\Omega_U$, both from \cite{wue26b}
(Proposition~\ref{prop:used}(ii),(i)); the third by an elementary estimate
(Lemma~\ref{lem:three}(c)), in which the blocks are normal and are
indefinite perturbations of the charges. The estimates at the different
vertices are combined with coefficients that increase along the arrows, and
$\gamma$ is the least increment; this is where the absence of oriented
cycles is used, and for a quiver with a loop no such bound holds
(Remark~\ref{rem:energy}). See Theorem~\ref{thm:energy} and
Proposition~\ref{prop:reparrows}.

\subsection{Cohomology that is not polystable}
Every nonzero representation $\mathcal H$ has a unique filtration
$0=\mathcal H^0\subset\mathcal H^1\subset\dots\subset\mathcal H^{\mathsf n}=\mathcal H$
by subrepresentations whose factors
$\operatorname{gr}^k=\mathcal H^k/\mathcal H^{k-1}$ are semistable with
$\arg Z(\operatorname{gr}^1)>\dots>\arg Z(\operatorname{gr}^{\mathsf n})$,
the Harder--Narasimhan filtration; we write
$\phi_k=\arg Z(\operatorname{gr}^k)$. The sum
$\sum_k|Z(\operatorname{gr}^k)|$ is the mass of $\mathcal H$ for the stability
condition.

\begin{theorem}[Harder--Narasimhan asymptotics]\label{thm:introH}
Let $\mathcal H$ be a nonzero representation of $Q$ with the phase flow from
an arbitrary initial metric, or the cohomology of a complex with a surjective
or an injective differential at one vertex, with the metrics $h_i(t)$ induced
by the phase flow of the complex. For all charges in the upper half-plane:
\begin{enumerate}
\item[(i)] for every nonzero vector $v$ of $\mathcal H_i$,
 $\lim_{t\to\infty}\frac1t\log h_i(t)(v,v)=-\phi_k$, where $k$ is the least
 index with $v\in\mathcal H^k_i$; so the Harder--Narasimhan filtration is the
 filtration of the spaces by the exponential rate of decay of the metric;
\item[(ii)] the eigenvalues of the argument of $z_i+\mu_i$, with $\mu$ the
 moment map of $\mathcal H$ for the metrics $h(t)$, converge to the $\phi_k$
 with the multiplicities $\dim\operatorname{gr}^k_i$, and
 $\sum_i\tr|z_i+\mu_i|$ converges to $\sum_k|Z(\operatorname{gr}^k)|$, which
 is its infimum over all metrics on $\mathcal H$; along the phase flow of a
 representation this sum, the mass, is nonincreasing;
\item[(iii)] the second fundamental forms of the subrepresentations
 $\mathcal H^k$ tend to zero exponentially, and the metric induced on
 $\operatorname{gr}^k$, multiplied by $e^{\phi_kt}$, has the properties
 stated in Theorem~\ref{thm:introB} and in Theorem~\ref{thm:introS} below
 for $\operatorname{gr}^k$ in the place of $\mathcal H$: its moment map
 tends to the stability parameter of $\operatorname{gr}^k$, and it converges
 exponentially to a harmonic metric if $\operatorname{gr}^k$ is polystable.
\end{enumerate}
\end{theorem}

\begin{theorem}[Semistable cohomology]\label{thm:introS}
Let $\mathcal H$ be as in Theorem~\ref{thm:introH} and semistable, with
$Z(\mathcal H)=|Z|e^{\ii\varphi}$. Then the metrics $e^{\varphi t}h(t)$ remain
at bounded distance, in the space of metrics, from a trajectory of the
gradient flow of the Kempf--Ness function of $\mathcal H$ for the
Riemannian metric $\sum_i\mathsf m_i\tr\xi_i^2$ with the weights
$\mathsf m_i=\Ima z_i/\sin^2\varphi$. Hence they are an asymptotic solution of
this gradient flow in the sense of \cite{hkkp23}, and by
\cite[Corollary~5.10 and Theorem~5.11]{hkkp23} their growth and decay are
given, up to bounded factors, by products of powers of iterated logarithms
of $t$ on the steps of the iterated weight filtration of $\mathcal H$, for
the weight function $\sum_i\mathsf m_i\dim\mathsf R_i$, which equals
$|Z(\mathsf R)|/\sin\varphi$ on the subquotients $\mathsf R$ of $\mathcal H$
of phase $\varphi$.
\end{theorem}

For a nonzero arrow $\CC\to\CC$ and two charges of equal argument, the
metrics grow and decay like powers of $t$ with exponents $|z_1|/|z_0+z_1|$
and $-|z_0|/|z_0+z_1|$. These exponents determine the ratio of the two
weights (Proposition~\ref{prop:weights}). The proof of
Theorem~\ref{thm:introS} compares the velocity of the flow with the gradient
of the Kempf--Ness function: the difference is quadratic in the shifted
moment map, hence integrable in time, and gradient flows of convex
functions on spaces of nonpositive curvature do not increase distances. The
proof of Theorem~\ref{thm:introH} uses the mass as a Lyapunov function; its
dissipation is bounded below by a multiple of the sum of the squared norms
of the commutators of the arrows with the phases. See Theorems~\ref{thm:comparison} and~\ref{thm:hn},
Corollary~\ref{cor:iterated} and
Propositions~\ref{prop:dissipation} and~\ref{prop:hnmass}.

\subsection{Kronecker quivers and the projective line}
Let $K_m$ be the generalized Kronecker quiver, with two vertices and $m$
arrows from the first to the second.

\begin{theorem}[Differentials at one vertex]\label{thm:introD}
Let $\mathsf E$ be a bounded complex of representations of $K_m$ whose
differentials vanish at one of the two vertices, with any Hermitian metric.
\begin{enumerate}
\item[(i)] $\mathsf E$ is an orthogonal direct sum of subcomplexes
 $\mathsf E_n$, one for each degree $n$, where $\mathsf E_n$ has cohomology
 $H^n(\mathsf E)$ in degree $n$ and, if $H^n(\mathsf E)\neq0$, is a shift of a
 complex with a surjective or with an injective differential at one vertex.
 The phase flow preserves this decomposition and is the product of the flows
 of the $\mathsf E_n$.
\item[(ii)] For all charges in the upper half-plane and every $n$, the
 blocks of the endomorphism $\Omega_{\mathsf E_n}$ of $\mathsf E_n$ approach
 those of $(-1)^n\Omega_{H^n(\mathsf E)}$, the endomorphism of the
 representation $H^n(\mathsf E)$ with its induced metrics, extended by zero;
 the trace-norm distance decays exponentially from a time that depends on
 the initial metric. On nonzero $H^n(\mathsf E)$ the induced metrics
 multiplied by $e^{-\pi nt}$ satisfy the conclusions of
 Theorem~\ref{thm:introH} and, in the semistable case, of
 Theorem~\ref{thm:introS}. On a Harder--Narasimhan quotient of phase
 $\phi$, the induced metric multiplied by $e^{(\phi-\pi n)t}$ converges
 exponentially to a harmonic metric if the quotient is polystable. The mass
 $\mass(h(t))$ converges to the sum over $n$ of the masses of the
 representations $H^n(\mathsf E)$ for the stability condition.
\item[(iii)] Minimal bounded complexes of projective representations of
 $K_m$, and minimal bounded complexes of injective representations, have
 differentials vanishing at one vertex. Hence (i) and (ii) hold for every
 object of the bounded derived category of $K_2$, that is, of the projective
 line, represented by such a complex with arbitrary metrics; the limit of
 the mass is the mass of the object for the stability condition with heart
 the category of representations.
\end{enumerate}
\end{theorem}

See Theorems~\ref{thm:decomposition}, \ref{thm:longtime} and~\ref{thm:P1}
and Proposition~\ref{prop:minimal}.

\subsection{Relation to earlier work and scope}
For $Q=K_m$ and a surjective differential at the first vertex, with $U\neq0$
and cohomology that is nonzero at both vertices,
Theorems~\ref{thm:introB}, \ref{thm:introA} and~\ref{thm:introC} are
\cite[Theorems~5.4, 5.7, 6.1 and~6.6]{wue26b}, with the same proofs, up to
two differences: here equal arguments of the two charges are included in
Theorem~\ref{thm:introB}, and the constant $c$ of
Theorem~\ref{thm:introC} is the one of \cite{wue26b} only after the
refinement of Remark~\ref{rem:energy}(2). Sections~\ref{sec:decay}
and~\ref{sec:convergence} follow \cite{wue26b}. What is added here is the
a priori bound on a general quiver (Section~\ref{sec:energy}), a reduction of
the Kempf--Ness theorem for real stability parameters to the rational case
of King (Lemma~\ref{lem:king}), the description of the flow for cohomology
that is not polystable (Section~\ref{sec:asymptotics}), and
Section~\ref{sec:kronecker}.
Theorem~\ref{thm:introR} is contained in \cite{wue26b} only for
representations of $K_m$ with both spaces nonzero and the charge of the
first vertex of larger argument than that of the second, where it follows
from the theorems on complexes by adding an exact summand.

Complexes with nonzero differentials at two or more vertices are open, in
particular complexes of representations of $K_m$ with differentials at both
vertices. Quivers with oriented cycles are excluded. For a cohomology
representation that is semistable and not polystable the induced metrics
have no limit after any normalization by a scalar function of the time
(Remark~\ref{rem:scope}); Theorem~\ref{thm:introS} determines them up to a
bounded factor, and whether this factor converges is open in general. It
converges in the example of Proposition~\ref{prop:weights} and for a
representation of $K_2$ for which the metrics are shown in
\cite[Proposition~6.9]{wue26b} to grow and decay like $t^{\pm1/2}$. The time $t_1$ of Theorem~\ref{thm:introA} is explicit
but very large (Remark~\ref{rem:decay}).

\subsection{Organization}
Section~\ref{sec:setting} fixes the setting and states the results used from
\cite{wue26a,wue26b,wue26c}. Section~\ref{sec:energy} proves
Theorem~\ref{thm:introC}, Section~\ref{sec:decay} Theorem~\ref{thm:introA},
and Section~\ref{sec:convergence} Theorems~\ref{thm:introR}
and~\ref{thm:introB}, including the Kempf--Ness theorem for real stability
parameters in the form needed. Section~\ref{sec:asymptotics} proves
Theorems~\ref{thm:introH} and~\ref{thm:introS}, and
Section~\ref{sec:kronecker} proves Theorem~\ref{thm:introD}.
